% RQ3 Theoretical Framework: Proposal Pipelines

We develop formal models of proposal pipeline dynamics and their effects on governance outcomes.

\subsection{Pipeline State Machine}

A proposal $p$ transitions through pipeline states:

\begin{equation}
p: S_0 \rightarrow S_1 \rightarrow ... \rightarrow S_{\text{terminal}}
\end{equation}

where terminal states are $\{\text{Passed}, \text{Failed}, \text{Expired}, \text{Abandoned}\}$.

\subsubsection{State Transition Probabilities}

At each stage, proposals may:
\begin{itemize}
    \item \textbf{Advance}: Pass gate, move to next stage
    \item \textbf{Fail}: Fail gate, terminate
    \item \textbf{Expire}: Exceed time limit, terminate
    \item \textbf{Abandon}: Proposer withdraws
\end{itemize}

\subsection{Temp-Check Model}

\subsubsection{Filtering Function}

Temp-checks filter proposals based on preliminary support signal:

\begin{equation}
\text{pass}_{\text{temp}}(p) = \mathbf{1}\left[\frac{\text{support votes}}{\text{total votes}} > \tau_{\text{temp}}\right]
\end{equation}

\subsubsection{Signal Quality}

Temp-check support is a noisy signal of eventual formal vote outcome:

\begin{equation}
\text{support}_{\text{temp}}(p) = \text{support}_{\text{true}}(p) + \epsilon
\end{equation}

where $\epsilon$ represents noise from:
\begin{itemize}
    \item Low temp-check turnout (sample variance)
    \item Different temp-check vs. formal vote populations
    \item Information revealed after temp-check
\end{itemize}

\subsubsection{Filtering Tradeoff}

Higher $\tau_{\text{temp}}$ reduces false positives (bad proposals reaching formal vote) but increases false negatives (good proposals filtered):

\begin{equation}
\text{FPR}(\tau) \downarrow \text{ as } \tau \uparrow; \quad \text{FNR}(\tau) \uparrow \text{ as } \tau \uparrow
\end{equation}

\subsection{Fast-Track Model}

\subsubsection{Qualification Criteria}

Proposals qualify for fast-track based on:

\begin{equation}
\text{fast-track}(p) = \mathbf{1}\left[\text{support}_{\text{early}}(p) > \tau_{\text{fast}}\right]
\end{equation}

Fast-tracked proposals skip intermediate stages, reducing time-to-decision.

\subsubsection{Fast-Track Risks}

Fast-tracking may:
\begin{itemize}
    \item Bypass deliberation on consequential issues
    \item Create capture opportunities (coordinate fast vote before opposition mobilizes)
    \item Reduce perceived legitimacy if overused
\end{itemize}

\subsection{Expiry Model}

\subsubsection{Expiry Dynamics}

Proposals expire if not resolved within window $\tau_{\text{expiry}}$:

\begin{equation}
\text{expired}(p, t) = \mathbf{1}\left[t - t_{\text{created}}(p) > \tau_{\text{expiry}}\right]
\end{equation}

\subsubsection{Abandonment vs. Expiry}

We distinguish:
\begin{itemize}
    \item \textbf{Expiry}: Proposal times out while still accumulating votes
    \item \textbf{Abandonment}: Proposer withdraws or stops advocating before expiry
\end{itemize}

Both indicate pipeline inefficiency but have different causes.

\subsection{Throughput-Quality Tradeoff}

\subsubsection{Throughput Metrics}

\begin{description}
    \item[Pass rate] Fraction of proposals that ultimately pass: $\frac{|\{p: \text{passed}(p)\}|}{|\text{proposals}|}$
    \item[Resolution rate] Fraction resolved (pass or fail) vs. abandoned/expired
    \item[Time-to-decision] Average time from creation to resolution
\end{description}

\subsubsection{Quality Proxies}

\begin{description}
    \item[Abandonment rate] Proposals abandoned mid-pipeline (suggests poor filtering)
    \item[Turnout] Participation in formal votes (higher suggests engagement)
    \item[Quorum rate] Fraction of formal votes achieving quorum
\end{description}

\subsection{Theoretical Predictions}

\subsubsection{Prediction 1: Temp-Check Filtering Curve}

Pass rate through formal voting increases with temp-check threshold (better filtering) but total proposal throughput decreases:

\begin{equation}
\text{pass rate} \uparrow \text{ with } \tau_{\text{temp}}; \quad \text{throughput} \downarrow \text{ with } \tau_{\text{temp}}
\end{equation}

\subsubsection{Prediction 2: Fast-Track Acceleration}

Fast-tracks reduce mean time-to-decision but increase variance (bimodal distribution):

\begin{equation}
\mathbb{E}[\text{TTD}] \downarrow; \quad \text{Var}[\text{TTD}] \uparrow
\end{equation}

\subsubsection{Prediction 3: Expiry-Abandonment Relationship}

Tight expiry windows increase abandonment rate for complex proposals:

\begin{equation}
\text{abandonment} \uparrow \text{ as } \tau_{\text{expiry}} \downarrow
\end{equation}
